\documentclass[11pt]{article}

\usepackage[margin=1in]{geometry}
\usepackage{amsmath}
\usepackage{booktabs,longtable,array}
\usepackage{hyperref}
\usepackage{xcolor}
\newcommand{\ch}[1]{\mbox{\texttt{#1}}}

\hypersetup{
  colorlinks=true,
  linkcolor=blue,
  citecolor=blue,
  urlcolor=blue
}

\title{Environment-Specific Prebiotic Reaction Networks for LIFEORIG}
\author{LIFEORIG working model specification}
\date{\today}

\begin{document}
\maketitle

\section{Purpose}

The previous \texttt{PREBIOTIC} reference file,
\texttt{prebiotic\_polymer\_autocatalytic\_network\_v2.txt}, is best treated
as a master scaffold, not as a single universal chemistry.  Early Earth
settings differ too strongly in energy source, phase structure, pH, pressure,
mineralogy, transport, and concentration mechanism for one reaction list to be
used without environmental filtering.  I therefore created three sibling
environment-specific networks:

\begin{center}
\begin{tabular}{lll}
\toprule
Environment & File & Number of reactions \\
\midrule
Hydrothermal vent &
\texttt{reference\_reactions/HYDRO\_VENT/hydro\_vent\_reaction\_network\_v1.txt} &
165 \\
Surface pond &
\texttt{reference\_reactions/SURFACE\_POND/surface\_pond\_reaction\_network\_v1.txt} &
175 \\
Volcanic rock &
\texttt{reference\_reactions/VOLCANIC\_ROCK\_ENV/volcanic\_rock\_reaction\_network\_v1.txt} &
170 \\
\bottomrule
\end{tabular}
\end{center}

These are not validated kinetic mechanisms.  They are literature-guided,
explicit, environment-conditioned hypothesis networks for simulation and
comparison.  Each line stores a reaction or lumped pathway, catalyst/control,
rate-law template, role, references, and confidence flag.

\section{Core Modeling Decision}

The correct comparison is not
\begin{quote}
one generic prebiotic network under different temperatures.
\end{quote}

The correct comparison is
\begin{quote}
a shared reaction library filtered into environment-specific active networks,
with different food sets, catalysts, kinetic modifiers, and loss terms.
\end{quote}

The overlap is chemically meaningful.  For example, \ch{HCN}, formamide,
glycolonitrile, aminoacetonitrile, glycine, phosphate species, and short
peptides can appear in more than one setting.  But the way they are produced,
stabilized, destroyed, and concentrated is environment-specific.

\section{Environment Distinctions}

\begin{longtable}{p{0.22\linewidth}p{0.24\linewidth}p{0.24\linewidth}p{0.22\linewidth}}
\toprule
Feature & Hydrothermal vent & Surface pond & Volcanic rock environment \\
\midrule
Dominant energy &
Redox disequilibrium, thermal gradients, pH gradients, fluid flow &
UV, evaporation, freezing/eutectic concentration, wet-dry cycles &
Volcanic heat, UV, gas deposition, wet-dry films, gas-rock redox \\
Dominant phase &
High-pressure aqueous pore/fluid system &
Shallow aqueous surface, air-water and ice-water interfaces &
Thin films on basalt, glass, ash, clay, sulfides, and pores \\
Primary catalysts &
\ch{FeS}, \ch{NiS}, magnetite, brucite, silica, pore barriers &
UV, sulfite/sulfide photoredox, clay, borate, phosphate minerals &
Basaltic glass, olivine alteration products, clay, \ch{FeS}, \ch{NiS},
magnetite, ash surfaces \\
Concentration mechanism &
Pore retention, adsorption, vent-ocean mixing, ice-interface HCN input &
Evaporation, freezing, surface enrichment, sediment adsorption &
Evaporating rock films, pore trapping, ash/glass adsorption, thermal cycling \\
Major losses &
Advection/dilution, hydrolysis, thermal degradation, oxidant mixing &
Photolysis, volatilization, washout, tar formation, precipitation &
Volatilization, runoff, thermal degradation, fouling, hydrolysis \\
\bottomrule
\end{longtable}

\section{Hydrothermal Vent Network}

The hydrothermal network emphasizes PHREEQC-controlled geochemistry and a
separate organic kinetic layer.  Acid-base equilibria, carbonate/sulfide/
phosphate speciation, mineral saturation, and precipitation should mostly be
handled by PHREEQC.  The organic layer then uses PHREEQC outputs as activities,
pH, ionic strength, redox state, and mineral surface availability.

The vent network has these main modules:
\begin{itemize}
\item geochemical speciation and precipitation: carbonate, sulfide, ammonia,
phosphate, \ch{FeS}, \ch{NiS}, magnetite, brucite, silica;
\item serpentinization and water-rock redox: \ch{H2} generation, methanation,
magnetite/brucite formation;
\item C1 fixation and sulfur carbonyl chemistry: formate, CO, COS,
formamide, thioformyl intermediates;
\item thioester and acetyl-phosphate activation: thioacetate,
acetyl phosphate, acyl transfer, hydrolysis losses;
\item central carboxylate and reductive amination chemistry: glyoxylate,
pyruvate, oxaloacetate, succinate, amino acids;
\item peptide, lipid, mineral-surface, transport, and degradation modules;
\item a warm HCN/ice-interface subnetwork from Das, Ghule, and Vanka, included
only as a boundary/interface branch rather than as the core vent chemistry.
\end{itemize}

Important examples include:
\begin{verbatim}
CO2 + H2 -> HCOO- + H+
CO2 + HS- + H+ -> COS + H2O
CO + CH3SH -> CH3COSR
CH3COSR + Pi -> acetyl_phosphate + RSH
glyoxylate + NH3 + H2 -> glycine + H2O
\end{verbatim}

\section{Surface Pond Network}

The surface pond network is the strongest home for HCN, UV, photoredox,
wet-dry, and freezing chemistry.  It should be used for shallow ponds,
lakeshores, tidal lagoons, ice/eutectic brines, and evaporating surface pools.
Unlike vents, it explicitly includes UV photolysis and atmospheric deposition as
first-class feedstock processes.

Its main modules are:
\begin{itemize}
\item atmospheric delivery of \ch{HCN}, \ch{H2CO}, \ch{NH3}, \ch{SO2},
\ch{H2S}, metals, phosphate, and meteoritic organics;
\item UV and sulfite/sulfide photoredox chemistry;
\item HCN oligomerization, formamide chemistry, nucleobase precursors;
\item formose chemistry with borate stabilization;
\item Strecker and reductive-amination amino-acid chemistry;
\item phosphate activation and wet-dry peptide/oligonucleotide condensation;
\item fatty-acid vesicle formation and wet-dry compartment cycling;
\item thermal aqueous HCN chemistry from Das, Ghule, and Vanka.
\end{itemize}

The Das--Ghule--Vanka paper is especially relevant here because it argues that
\ch{HCN} and water alone can generate key RNA/protein precursors under
approximately 80--100 \(^\circ\)C thermal aqueous conditions, with water or
ammonia acting as proton shuttles and without requiring metal catalysts.
Reactions added from that subnetwork include:
\begin{verbatim}
HCN + H2O -> formamide
formamide + H2O -> formic_acid + NH3
formic_acid -> CO + H2O
HCN + H2O -> formaldimine
formaldimine + HCN -> aminoacetonitrile
aminoacetonitrile + 2H2O -> glycine + NH3
urea -> cyanamide + H2O
formamide + glycolonitrile -> amino_oxazole_derivative + H2O
\end{verbatim}

\section{Volcanic Rock Network}

The volcanic rock network is not simply a dry version of the surface pond.  It
is centered on rock films, ash beds, basaltic glass, volcanic gases, fumarolic
inputs, heat pulses, and mineral surfaces.  It includes some HCN and wet-dry
chemistry, but the boundary conditions are gas-rock-film and thermal cycling,
not open pond photochemistry alone.

Its main modules are:
\begin{itemize}
\item volcanic gas input: \ch{CO2}, CO, \ch{H2}, \ch{H2S}, \ch{SO2}, COS,
\ch{HCN}, \ch{H2CO}, \ch{NH3}, NOx;
\item basaltic glass, olivine, ash, clay, apatite, carbonate, and sulfide
weathering;
\item sulfur-carbonyl and Fe/Ni sulfide surface chemistry;
\item C1/C2 reduction and carboxylate formation in thin films;
\item HCN/nitrile, formose, Strecker, and thermal HCN-water branches;
\item dry heat, wet-dry, mineral-surface peptide and nucleotide polymerization;
\item fatty-acid film rehydration, vesicle formation, and pore transport/loss.
\end{itemize}

Volcanic rock settings therefore bridge atmospheric/surface chemistry and
mineral redox chemistry.  They can receive atmospheric HCN and formaldehyde,
but also offer \ch{FeS}/\ch{NiS}, magnetite, basalt glass, phosphate minerals,
and dry thermal activation.

\section{How To Use These Networks}

For computation, the recommended path is:
\begin{enumerate}
\item Treat each file as an environment-specific candidate network.
\item Parse reaction IDs, modules, reactants, products, controls, references,
and confidence flags.
\item Build a food set from the environment solver: atmosphere, hydrology,
PHREEQC, rock weathering, gas fluxes, UV, and mineral availability.
\item Activate reactions only when required feedstocks and environmental
controls exist.
\item Use PHREEQC for geochemical equilibria and mineral saturation where
flagged; use custom kinetics for organic nonequilibrium pathways.
\item Compare networks by reachable products, autocatalytic closure,
feedstock cost, robustness to losses, and overlap with extant metabolic motifs.
\end{enumerate}

\section{Immediate Caveats}

These networks are intentionally large, but they should not be mistaken for
rate-validated mechanisms.  The confidence flags are important.  High-confidence
entries are generally speciation, hydrolysis, adsorption, precipitation, or
well-established feedstock transformations.  Medium-confidence entries are
plausible prebiotic transformations with literature support but uncertain
environmental yield.  Low-confidence entries are useful hypothesis links,
coarse-grained pathways, or autocatalytic/protocell test rules.

\section{References}

\begin{thebibliography}{99}

\bibitem{miller1953}
S. L. Miller,
A production of amino acids under possible primitive Earth conditions,
\emph{Science}, 1953, 117, 528--529.

\bibitem{ferris1984}
J. P. Ferris and W. J. Hagan,
HCN and chemical evolution: the possible role of cyano compounds in prebiotic synthesis,
\emph{Tetrahedron}, 1984, 40, 1093--1120.

\bibitem{powner2009}
M. W. Powner, B. Gerland, and J. D. Sutherland,
Synthesis of activated pyrimidine ribonucleotides in prebiotically plausible conditions,
\emph{Nature}, 2009, 459, 239--242.

\bibitem{ritson2012}
D. Ritson and J. D. Sutherland,
Prebiotic synthesis of simple sugars by photoredox systems chemistry,
\emph{Nature Chemistry}, 2012, 4, 895--899.

\bibitem{patel2015}
B. H. Patel, C. Percivalle, D. J. Ritson, C. D. Duffy, and J. D. Sutherland,
Common origins of RNA, protein and lipid precursors in a cyanosulfidic protometabolism,
\emph{Nature Chemistry}, 2015, 7, 301--307.

\bibitem{das2019}
T. Das, S. Ghule, and K. Vanka,
Insights into the origin of life: did it begin from HCN and H2O?,
\emph{ACS Central Science}, 2019, 5, 1532--1540.
\url{https://doi.org/10.1021/acscentsci.9b00520}

\bibitem{fabian2014}
B. Fabian, M. Szori, and P. Jedlovszky,
Floating patches of HCN at the surface of their aqueous solutions,
\emph{Journal of Physical Chemistry C}, 2014, 118, 21469--21482.

\bibitem{szori2014}
M. Szori and P. Jedlovszky,
Adsorption of HCN at the surface of ice,
\emph{Journal of Physical Chemistry C}, 2014, 118, 3599--3609.

\bibitem{martinRussell2007}
W. Martin and M. J. Russell,
On the origin of biochemistry at an alkaline hydrothermal vent,
\emph{Philosophical Transactions of the Royal Society B}, 2007, 362, 1887--1925.

\bibitem{sojo2016}
V. Sojo, B. Herschy, A. Whicher, E. Camprubi, and N. Lane,
The origin of life in alkaline hydrothermal vents,
\emph{Astrobiology}, 2016, 16, 181--197.

\bibitem{herschy2014}
B. Herschy et al.,
An origin-of-life reactor to simulate alkaline hydrothermal vents,
\emph{Journal of Molecular Evolution}, 2014, 79, 213--227.

\bibitem{huber1997}
C. Huber and G. Wachtershauser,
Activated acetic acid by carbon fixation on (Fe,Ni)S under primordial conditions,
\emph{Science}, 1997, 276, 245--247.

\bibitem{cody2000}
G. D. Cody et al.,
Primordial carbonylated iron-sulfur compounds and the synthesis of pyruvate,
\emph{Science}, 2000, 289, 1337--1340.

\bibitem{whicher2018}
A. Whicher et al.,
Acetyl phosphate as a primordial energy currency at the origin of life,
\emph{Origins of Life and Evolution of Biospheres}, 2018, 48, 159--179.

\bibitem{mccollom2007}
T. M. McCollom and J. S. Seewald,
Abiotic synthesis of organic compounds in deep-sea hydrothermal environments,
\emph{Chemical Reviews}, 2007, 107, 382--401.

\bibitem{deamer2017}
D. W. Deamer,
The role of lipid membranes in life's origin,
\emph{Life}, 2017, 7, 5.

\bibitem{damer2020}
B. Damer and D. Deamer,
The hot spring hypothesis for an origin of life,
\emph{Astrobiology}, 2020, 20, 429--452.

\bibitem{ferris1996}
J. P. Ferris, A. R. Hill, R. Liu, and L. E. Orgel,
Synthesis of long prebiotic oligomers on mineral surfaces,
\emph{Nature}, 1996, 381, 59--61.

\bibitem{ricardo2004}
A. Ricardo, M. A. Carrigan, A. N. Olcott, and S. A. Benner,
Borate minerals stabilize ribose,
\emph{Science}, 2004, 303, 196.

\bibitem{parkhurst2013}
D. L. Parkhurst and C. A. J. Appelo,
Description of input and examples for PHREEQC version 3,
U.S. Geological Survey Techniques and Methods, 2013.
\url{https://water.usgs.gov/water-resources/software/PHREEQC/}

\end{thebibliography}

\end{document}
